\emph{Setting.} The frame, $X(x)$, $I_{\Lambda}$ and the scale $\sigma = \ip{X(x)}{I_{\Lambda}}$ are those of Lemma \claimref{mass.de-sitter-scale-frozen-point}; $W$, $B_{\Lambda} = e_{0}W$, its plane, the unit boost, $R_{h}$, $\kappa_{\Lambda}$ and the bifurcation sphere $S_{\Lambda}$ with the horizon at areal radius $R = r/\sigma = R_{h}$ on the slice $x = re_{1}$ are those of Proposition \claimref{synthesis.de-sitter-static-boost-frozen-point}, read over $K$. Nothing in either uses a square root except the clauses analysed here: the identities of their clauses (i) are polynomial in $\Lambda$ and so hold over any field of characteristic $0$ (division by $2$, $3$ and $6$ only). For a $K$-vector space with an orthogonal basis $b_{1}, \dots, b_{n}$ the Clifford algebra $\Cl_{K}$ has the basis of products $b_{S}$, $S \subseteq \{1, \dots, n\}$, with $b_{i}b_{j} = -b_{j}b_{i}$ ($i \ne j$) and $b_{i}^{2} = Q(b_{i})$; an isometry carries such a basis to one with the same squares and so induces an algebra isomorphism.

\emph{(i).} $Q(I_{\Lambda}) = 2\cdot\tfrac{\Lambda}{6}\ip{\nO}{\nInf} = d$, $Q(W) = -d$, $\ip{W}{I_{\Lambda}} = \tfrac{\Lambda}{6} - \tfrac{\Lambda}{6} = 0$, and $B_{\Lambda}^{2} = e_{0}We_{0}W = -e_{0}^{2}W^{2} = d$, as in the Proposition. $\ip{X(x)}{I_{\Lambda}} = \tfrac{\Lambda}{6}\ip{\nO}{\nO} + \ip{\nO}{\nInf} - \tfrac12 Q(x)\tfrac{\Lambda}{6} = 1 - \Lambda Q(x)/12$. The vectors $e_{0}, \dots, e_{3}, W$ are orthogonal to $I_{\Lambda}$ and to each other, with the squares listed, and together with $I_{\Lambda}$ they span the six-dimensional space ($\nInf = \tfrac12(I_{\Lambda} + W)$, $\nO = \tfrac{3}{\Lambda}(I_{\Lambda} - W)$), so they are a basis of $I_{\Lambda}^{\perp}$. With $e_{0}^{2} = 1$ and $We_{0} = -e_{0}W$: $B_{\Lambda}e_{0} - e_{0}B_{\Lambda} = -W - W$, and $B_{\Lambda}W - WB_{\Lambda} = e_{0}W^{2} - We_{0}W = -d\,e_{0} - d\,e_{0}$; a vector $y$ orthogonal to $e_{0}$ and $W$ anticommutes with both, so commutes with $B_{\Lambda}$. The matrix of $\tfrac12\mathrm{ad}\,B_{\Lambda}$ on $(e_{0}, W)$ is $\bigl(\begin{smallmatrix}0 & -d\\ -1 & 0\end{smallmatrix}\bigr)$, with characteristic polynomial $t^{2} - d$. In the five-dimensional $\Cl_{K}(I_{\Lambda}^{\perp})$, a basis vector $b_{i}$ satisfies $b_{i}b_{S} = (-1)^{|S|}b_{S}b_{i}$ for $i \notin S$ and $(-1)^{|S|-1}b_{S}b_{i}$ for $i \in S$, so $b_{S}$ is central exactly when $S = \emptyset$ or $S$ is the whole basis ($|S| = 5$); since conjugation by $b_{i}$ is diagonal on the blades, the centre is $\mathrm{span}\{1, \omega\}$. Reversing the order of five factors costs $10$ transpositions, so $\omega^{2} = \prod_{i}Q(b_{i}) = 1\cdot(-1)^{3}\cdot(-d) = d$, and with the form reversed $\omega'^{2} = (-1)\cdot 1\cdot d = -d$. The check verifies all of this exactly over $\Q$ at eleven values of $\Lambda$ of both signs and several square classes.

\emph{(iii), hence (a) implies (b) to (j).} Let $\Lambda = 3H^{2}$, $H \in K$, $H \neq 0$. $Q(W/H) = -1$, so $v, v^{*}$ have $Q = \tfrac14(1 - 1) = 0$, $vv^{*} + v^{*}v = 2\ip{v}{v^{*}} = \tfrac12(1 + 1) = 1$ and $v^{*}v - vv^{*} = \tfrac12(e_{0}e_{W} - e_{W}e_{0}) = e_{0}e_{W} = u$ with $e_{W} = W/H$; $u = B_{\Lambda}/H$ has $u^{2} = d/H^{2} = 1$. By (i), $\tfrac12\mathrm{ad}\,B_{\Lambda}$ sends $e_{0} + e_{W} \mapsto -He_{W} - He_{0}$ and $e_{0} - e_{W} \mapsto -He_{W} + He_{0}$, so $v \mapsto -Hv$, $v^{*} \mapsto Hv^{*}$. $R_{h} = 1/H$ has $R_{h}^{2} = 1/d = 3/\Lambda$ and $\kappa = H$ has $\kappa^{2} = d$. For $x = (2/H)e_{1}$, $Q(x) = -4/H^{2} = -12/\Lambda$, so $x \in S_{\Lambda}$, $\sigma = 1 + 4\Lambda/(12H^{2}) = 2$ and $R = (2/H)/2 = R_{h}$, the horizon point of the Proposition. The basis $e_{0}, \dots, e_{3}, e_{W}$ is orthonormal with squares $\langle 1, -1, -1, -1, -1\rangle$: (h), and so (i) and (j) by the induced isomorphisms (the reversed forms are isometric as well). Directly: $(\omega/H)^{2} = 1$, so $\tfrac12(1 \pm \omega/H)$ are central idempotents, and $(\omega'/H)^{2} = -1$.

\emph{(b) to (j) each imply (a).} (b): $Q(ae_{0} + bW) = a^{2} - db^{2}$; if this vanishes with $(a, b) \neq 0$ then $b \neq 0$ and $d = (a/b)^{2}$. (c): $\Pi_{\Lambda}$ is nondegenerate ($d \neq 0$), so a vector is $x = x_{\Pi} + x_{\perp}$; $\tfrac12\mathrm{ad}\,B_{\Lambda}$ preserves $\Pi_{\Lambda}$ and kills $\Pi_{\Lambda}^{\perp}$, so $tx_{\perp} = 0$, $x_{\perp} = 0$, and $x_{\Pi}$ is an eigenvector on $\Pi_{\Lambda}$, so $t^{2} = d$. (d): $(cB_{\Lambda})^{2} = c^{2}d = 1$ gives $d = (1/c)^{2}$. (e): $Q(cW) = -c^{2}d = -1$, the same. (f): $R_{h}^{2} = 1/d$ gives $d = (1/R_{h})^{2}$; $\kappa^{2} = d$ is (a). (g): $re_{1} \in S_{\Lambda}$ means $-r^{2} = -12/\Lambda = -4/d$, so $d = (2/r)^{2}$. (h) implies (i) by the induced isomorphism. (i): an algebra isomorphism carries centre onto centre; the centre of $\Cl_{K}(1,4)$ is $\mathrm{span}\{1, \omega_{0}\}$ with $\omega_{0}^{2} = 1$ (the computation of (i) with $d = 1$), which contains the idempotents $\tfrac12(1 \pm \omega_{0})$; an idempotent $a + b\omega$ of $\mathrm{span}\{1, \omega\}$ has $a^{2} + db^{2} = a$ and $2ab = b$, so $b \neq 0$ forces $a = \tfrac12$ and $db^{2} = \tfrac14$, $d = (1/2b)^{2}$. (j): the centre of $\Cl_{K}(4,1)$ is $\mathrm{span}\{1, \omega_{1}\}$ with $\omega_{1}^{2} = -1$; an isomorphism sends $\omega_{1}$ to some $a + b\omega'$ with $(a + b\omega')^{2} = a^{2} - db^{2} + 2ab\,\omega' = -1$, so $ab = 0$. If $a = 0$, $d = (1/b)^{2}$. If $b = 0$, $a^{2} = -1$, so $-1 \in K^{2}$; then $\omega_{1}/a$ has square $1$ and $\mathrm{span}\{1, \omega_{1}\}$ contains the idempotent $\tfrac12(1 + \omega_{1}/a)$, so $\mathrm{span}\{1, \omega'\}$ contains one, which by the computation just made for (i) (with $-d$ in place of $d$) gives $-d \in K^{2}$, and with $-1 \in K^{2}$, $d \in K^{2}$. This proves (ii). The check verifies the reductions over $\Q$ at the eight sampled $\Lambda$ with $d \notin \Q^{2}$, where it decides squareness exactly, computes the two-dimensional centre by exact linear algebra at $\Lambda = 6$ and $7/3$, and verifies the witnesses of (iii) at five rational $H$ and, over the quadratic fields $\mathrm{span}\{1, u\}$ ($u$ the product of two extra generators of squares $1$ and $-\delta$, central on the blades of $\Cl(2,4)$, $u^{2} = \delta$), at $\Lambda = 6, 7/3, 1, 2/3$ with $H$ a rational multiple of $u$.

\emph{(iv).} $7/9$ and $2$ are not rational squares, so (b) to (j) fail by (ii). A rational point of $S_{7/3}$ has $x_{1}^{2} + x_{2}^{2} + x_{3}^{2} = 36/7$, and $y = 7x/6$ then has $y_{1}^{2} + y_{2}^{2} + y_{3}^{2} = 7$; clearing denominators, $A^{2} + B^{2} + C^{2} = 7D^{2}$ with $\gcd(A, B, C, D) = 1$. Squares modulo $8$ are $0, 1, 4$, so a sum of three is never $7$ modulo $8$, and $D$ is even; then $A^{2} + B^{2} + C^{2} \equiv 0 \pmod 4$, and since squares modulo $4$ are $0, 1$, all of $A, B, C$ are even, against the gcd. So $S_{7/3}$ has no rational point. At $\Lambda = 6$, $\mathrm{span}\{1, \omega\}$ with $\omega^{2} = 2$ is a field ($(a + b\omega)(a - b\omega) = a^{2} - 2b^{2} \neq 0$ for $(a, b) \neq 0$), and an isomorphism of $\mathrm{span}\{1, \omega'\}$, $\omega'^{2} = -2$, with $\mathrm{span}\{1, \omega_{1}\}$, $\omega_{1}^{2} = -1$, would need $a^{2} - 2b^{2} = -1$ and $ab = 0$, that is $a^{2} = -1$ or $b^{2} = \tfrac12$, neither possible in $\Q$; the two fields are those of discriminants $-8$ and $-4$. $Q(e_{1} + e_{2}) = -2 = -12/6$ puts $e_{1} + e_{2}$ on $S_{6}$, and $2 = 1^{2} + 1^{2}$ gives (v)(b). The check runs the residue tables modulo $8$ and $4$ completely, a search with no solution of $A^{2} + B^{2} + C^{2} = 7D^{2}$ for $D \le 15$, and the computations at $\Lambda = 6$.

\emph{(v)(a).} $x \in S_{\Lambda}$ means $x_{0} = 0$ and $x_{1}^{2} + x_{2}^{2} + x_{3}^{2} = 4/d$, equivalently $\sum_{i}(dx_{i}/2)^{2} = d$. If $d = H^{2}$, $x = (2/H)e_{1}$. $\Lambda = 6$ and $7/3$ are (iv).

\emph{(v)(b).} If $d = a^{2} + b^{2}$, let $f_{0}, \dots, f_{4}$ be the standard basis of squares $\langle 1, -1, -1, -1, -1\rangle$ and $T e_{0} = \tfrac{1+d}{2}f_{0} + \tfrac{1-d}{2}f_{1}$, $T e_{1} = \tfrac{1-d}{2}f_{0} + \tfrac{1+d}{2}f_{1}$, $T e_{2} = af_{2} + bf_{3}$, $T e_{3} = -bf_{2} + af_{3}$, $T W = d\,f_{4}$. Then $Q(Te_{0}) = \tfrac14((1+d)^{2} - (1-d)^{2}) = d$, $Q(Te_{1}) = -d$, $Q(Te_{2}) = Q(Te_{3}) = -(a^{2} + b^{2}) = -d$, $Q(TW) = -d^{2} = d\,Q(W)$, and the images are orthogonal, so $Q(Tx) = d\,Q(x)$; $T$ is invertible ($\det T = d^{3}$). Conjugation $g \mapsto TgT^{-1}$ preserves $Q$ up to the factor, which cancels, and the determinant, so it carries $SO(I_{\Lambda}^{\perp})$ onto $SO_{K}(1,4)$, and the induced map on bivectors carries the bivector Lie algebra onto $\mathfrak{so}_{K}(1,4)$. Conversely let $Q(Tx) = c\,Q(x)$. The images $Tb_{i}$ of the orthogonal basis are orthogonal, so their product is their outer product, $\det T$ times $f_{0}\cdots f_{4}$, whose square is $1$; computed from the factors, its square is $\prod_{i}c\,Q(b_{i}) = c^{5}d$. So $c^{5}d = (\det T)^{2}$, $cd \in K^{2}$, and $c = ds^{2}$ for some $s \in K$; replacing $T$ by $T/s$, $d\,Q$ on $I_{\Lambda}^{\perp}$, which is $\langle d, -d, -d, -d, -d^{2}\rangle$, is isometric to $\langle 1, -1, -1, -1, -1\rangle$. By the map above with $(a, b) = (1, 0)$ restricted to the first two vectors, $\langle d, -d\rangle \cong \langle 1, -1\rangle$, and $\langle -d^{2}\rangle \cong \langle -1\rangle$ by $W \mapsto W/d$. So $\langle 1, -1, -1\rangle \perp \langle -d, -d\rangle \cong \langle 1, -1, -1\rangle \perp \langle -1, -1\rangle$. One-dimensional cancellation, proved by a reflection: if $\phi$ is an isometry of $\langle c\rangle \perp U$ onto $\langle c\rangle \perp U'$, $c \neq 0$, with $y$ spanning the first summand of the target and $z = \phi(y_{0})$ the image of the generator $y_{0}$ of the source's, then $Q(z) = Q(y) = c$ and $Q(z - y) + Q(z + y) = 4c \neq 0$, so $n = z - y$ or $n = z + y$ is not null, and the versor $x \mapsto -nxn^{-1}$, the reflection $x \mapsto x - 2\ip{x}{n}n/Q(n)$, sends $z$ to $y$ or to $-y$ (since $2\ip{z}{z \mp y} = Q(z \mp y)$); composing, $y_{0} \mapsto \pm y$, so $\phi$ followed by the reflection carries $U = y_{0}^{\perp}$ isometrically onto $U' = y^{\perp}$. Cancelling $\langle 1\rangle$, $\langle -1\rangle$, $\langle -1\rangle$ in turn, $\langle -d, -d\rangle \cong \langle -1, -1\rangle$, so $\langle d, d\rangle \cong \langle 1, 1\rangle$ and $\langle 1, 1\rangle$ represents $d$: $d = a^{2} + b^{2}$. If $d = H^{2}$, $d = H^{2} + 0^{2}$; $\Lambda = 6$ is (iv). For $\Lambda = 9$, $3 = (a/c)^{2} + (b/c)^{2}$ would give $a^{2} + b^{2} = 3c^{2}$ with $\gcd(a, b, c) = 1$; squares modulo $3$ are $0, 1$, so $3 \mid a$ and $3 \mid b$, then $9 \mid 3c^{2}$ and $3 \mid c$, against the gcd. The check verifies $Q(Tx) = d\,Q(x)$ on all pairs of basis vectors and $\det T$ through the pseudoscalar at six sums of two squares, the discriminant step, the reflection on three pairs of vectors, the residue table modulo $3$, and a search with no solution of $a^{2} + b^{2} = 3c^{2}$ for $c \le 24$.

\emph{(v)(c).} The centre $\mathrm{span}\{1, \omega'\}$, $\omega'^{2} = -d$, is a field exactly when it has no idempotent other than $0, 1$, which by the computation in (i) is $-d \notin K^{2}$. With $d = H^{2}$ this is $-1 \notin K^{2}$, true in an ordered field, where squares are not negative. Over $\mathrm{span}\{1, u\}$, $u^{2} = -1$, $u$ commutes with $\omega'$, so $(u\omega'/H)^{2} = u^{2}\omega'^{2}/H^{2} = 1$ and $\tfrac12(1 + u\omega'/H)$ is a central idempotent; the check verifies this at two $H$. With (ii)(i) and (j) this gives the stated reading.

\emph{(v)(d).} $Q(W/c) = -d/c^{2}$; for $c^{2} = |\Lambda|$ this is $-\Lambda/(3|\Lambda|) = -\mathrm{sgn}(\Lambda)/3$, so the rescaled complement is $\Cl_{K}(1,3) \perp \langle \pm\tfrac13\rangle$ with $+$ for $\Lambda < 0$. Its central element has square $\mp\tfrac13$ by the computation of (i), against $\mp 1$ for $\Cl_{K}(2,3)$ and $\Cl_{K}(1,4)$. Upper sign: an isomorphism of centres sends the unit of square $-1$ to $a + b\omega$ with $a^{2} - b^{2}/3 = -1$ and $ab = 0$; $a = 0$ gives $3 = b^{2}$, and $b = 0$ gives $-1 \in K^{2}$, then both centres split, $-\tfrac13 \in K^{2}$ and $3 = (-1)/(-\tfrac13) \in K^{2}$. Lower sign: the centre of $\Cl_{K}(1,4)$ splits, so $\tfrac13 \in K^{2}$ by the idempotent computation. Conversely, if $3 = t^{2}$, then $t\,W/c$ has square $\pm 1$ and the forms are isometric. With $H \in K$, $|\Lambda| = 3H^{2}$ is a square if and only if $3$ is. The check verifies the squares at five $\Lambda$ with $|\Lambda|$ a square, the vector of square $\pm1$ over $\mathrm{span}\{1, u\}$, $u^{2} = 3$, and $Q(W/H) = \mp 1$ at five $\Lambda = \pm 3H^{2}$.

\emph{(v)(e).} $(cB_{\Lambda})^{2} = c^{2}d = -1$ and $Q(cW) = -c^{2}d = 1$ each give $-d = (1/c)^{2}$, and conversely $c = 1/H'$. $I_{\Lambda}^{\perp} \cong \langle 1, -1, -1, -1, 1\rangle$ implies $\Cl_{K}(I_{\Lambda}^{\perp}) \cong \Cl_{K}(2,3)$, whose centre has unit $\omega_{2}$ with $\omega_{2}^{2} = -1$; as in (j), an isomorphism of centres gives $a^{2} + db^{2} = -1$ with $ab = 0$, so $-d = (1/b)^{2}$, or $a^{2} = -1$, both centres split and $d, -1 \in K^{2}$, so $-d \in K^{2}$. With $H'$, $W/H'$ completes the basis to $\langle 1, -1, -1, -1, 1\rangle$. Over an ordered $K$ with $d < 0$, $Q(ae_{0} + bW) = a^{2} + |d|b^{2} > 0$ unless $a = b = 0$. The check verifies the witnesses at two $H'$ and the definiteness on a grid.

\emph{(v)(f).} $\lambda s^{4} = 1$ says $(s^{2})^{2} = 1/\lambda$, so $s^{2} \in K$ gives $\lambda = (1/s^{2})^{2}$, and conversely $s^{2} = 1/\sqrt{\lambda}$. With $\Lambda = 3H^{2}$, $\lambda = H^{2}/2$ is a square exactly when $\tfrac12$ is, that is $2 \in K^{2}$. The check verifies that $\lambda$ is not a rational square at three rational $H$, that $s^{2} = u/H$ solves $\lambda s^{4} = 1$ over $\mathrm{span}\{1, u\}$, $u^{2} = 2$, and that the stored sample points $\Lambda = 6$ and $\tfrac32$ have $\sqrt{\lambda} = 1, \tfrac12$ rational and $\Lambda/3$ not a square.
